% =============================================================================
% Capítulo 5 – Modelagem e Diagramas do SIGEA-GTT
% =============================================================================
\chapter{Modelagem e Diagramas}
\label{cap_modelagem}

Este capítulo apresenta os principais artefatos de modelagem do SIGEA-GTT,
produzidos em conformidade com as notações UML 2.5 \cite{booch2005} e com o
Modelo de Arquitetura 4+1 de \citeonline{kruchten1995}. Os diagramas foram
elaborados com a biblioteca TikZ e expressam as decisões arquiteturais que
fundamentam a implementação descrita nos capítulos anteriores.

% ---------------------------------------------------------------------------
\section{Diagrama de Entidade-Relacionamento (DER)}
\label{sec_der}
% ---------------------------------------------------------------------------

O Diagrama de Entidade-Relacionamento (DER) da \autoref{fig_der} modela as
entidades persistidas no banco de dados PostgreSQL 16 e suas associações.
O esquema foi versionado por Flyway e gerenciado pelo Hibernate (JPA) com
estratégia de DDL \textit{validate}. As principais entidades são:
\texttt{users}, \texttt{turmas}, \texttt{atividades}, \texttt{submissoes},
\texttt{modulos\_gtt}, \texttt{gatilhos\_gtt} e \texttt{gravidades\_merp}.

\begin{figure}[htb]
\centering
\caption{Diagrama de Entidade-Relacionamento (DER) simplificado do SIGEA-GTT.}
\label{fig_der}
\begin{tikzpicture}[
  >=Stealth,
  entity/.style={
    rectangle, draw=blue!70!black, fill=blue!8, thick,
    minimum width=3.0cm, minimum height=0.8cm,
    font=\small\bfseries, rounded corners=2pt
  },
  rela/.style={
    diamond, draw=orange!80!black, fill=yellow!10, thick,
    minimum width=2.2cm, minimum height=0.7cm,
    font=\tiny\bfseries, aspect=2.5
  },
  lbl/.style={font=\tiny, midway}
]

% -- Linha superior de entidades --
\node[entity] (users)     at ( 0,  0) {users};
\node[entity] (turmas)    at ( 5,  0) {turmas};
\node[entity] (atividades)at (10,  0) {atividades};

% -- Linha inferior de entidades --
\node[entity] (submissoes)at (10, -4) {submissoes};
\node[entity] (modulos)   at ( 0, -4) {modulos\_gtt};
\node[entity] (gatilhos)  at ( 5, -4) {gatilhos\_gtt};
\node[entity] (merp)      at ( 0, -8) {gravidades\_merp};

% -- Relacionamentos --
\node[rela] (r1) at (2.5,  0) {leciona};
\node[rela] (r2) at (7.5,  0) {possui};
\node[rela] (r3) at (5,   -2) {matricula};
\node[rela] (r4) at (10, -2)  {gera};
\node[rela] (r5) at (2.5, -4) {contem};
\node[rela] (r6) at (7.5, -4) {referencia};
\node[rela] (r7) at ( 0,  -6) {classifica};

% -- Conexões --
\draw[-] (users) -- node[lbl, above]{1} (r1);
\draw[->](r1)    -- node[lbl, above]{N} (turmas);

\draw[-] (turmas)    -- node[lbl, above]{1} (r2);
\draw[->](r2)        -- node[lbl, above]{N} (atividades);

\draw[-] (turmas)    -- node[lbl, left]{N}  (r3);
\draw[->](r3)        -- node[lbl, right]{M} (users);

\draw[-] (atividades)-- node[lbl, right]{1} (r4);
\draw[->](r4)        -- node[lbl, right]{N} (submissoes);

\draw[->](users)     -- node[lbl, left]{1:N} (submissoes);

\draw[-] (modulos)   -- node[lbl, above]{1} (r5);
\draw[->](r5)        -- node[lbl, above]{N} (gatilhos);

\draw[-] (gatilhos)  -- node[lbl, above]{N} (r6);
\draw[->](r6)        -- node[lbl, above]{M} (submissoes);

\draw[-] (merp)      -- node[lbl, right]{1} (r7);
\draw[->](r7)        -- node[lbl, right]{N} (gatilhos);

\end{tikzpicture}
\fonte{Elaborado pelo autor com TikZ (2026).}
\end{figure}

% ---------------------------------------------------------------------------
\section{Diagrama de Casos de Uso}
\label{sec_casos_uso}
% ---------------------------------------------------------------------------

O Diagrama de Casos de Uso da \autoref{fig_usecase} define o escopo
funcional do sistema a partir da perspectiva dos três atores descritos no
\autoref{cap_requisitos}. Os casos de uso foram derivados diretamente dos
Requisitos Funcionais RF01--RF25 e organizados em agrupamentos por perfil
de acesso.

\begin{figure}[p]
\centering
\caption{Diagrama de Casos de Uso do SIGEA-GTT (UML 2.5).}
\label{fig_usecase}
\begin{tikzpicture}[
  >=Stealth,
  uc/.style={
    ellipse, draw=blue!70!black, fill=blue!5, thick,
    font=\tiny, align=center, minimum width=3.2cm, minimum height=0.8cm
  },
  bdr/.style={
    rectangle, draw=blue!40, dashed, thick, rounded corners=4pt,
    fill=blue!2
  },
  act/.style={font=\small\bfseries}
]

% Sistema Boundary
\node[bdr, minimum width=9cm, minimum height=14cm]
     (sys) at (5, -5) {};
\node[above, font=\small\bfseries, blue!70!black] at (sys.north) {SIGEA-GTT};

% Atores
\node[act] (admin) at (-2,  -1) {ADMIN};
\node[act] (prof)  at (-2,  -6) {PROFESSOR};
\node[act] (aluno) at (12,  -7) {ALUNO};

% -- Casos de uso comuns --
\node[uc] (login)   at (5,  0) {Login / JWT\\(RF01, RF02)};
\node[uc] (senha)   at (5, -1.5) {Alterar Senha\\(RF03)};
\node[uc] (catalogo)at (5, -3)   {Catálogo GTT\\e MERP (RF08, RF09)};

% -- Admin --
\node[uc] (uc_users)   at (3, -4.5)  {Gerenciar\\Usuários (RF04, RF05)};
\node[uc] (uc_modgtt)  at (3, -6.0)  {CRUD Módulos\\GTT (RF06)};
\node[uc] (uc_gat)     at (3, -7.5)  {CRUD Gatilhos\\GTT (RF07)};
\node[uc] (uc_turmas)  at (3, -9.0)  {Gerenciar Turmas\\(RF10--RF13)};
\node[uc] (uc_dash_a)  at (3, -10.5) {Dashboard\\Analítico (RF25)};

% -- Professor --
\node[uc] (uc_atv)     at (7, -4.5)  {Criar Atividade\\(RF14, RF15)};
\node[uc] (uc_corr)    at (7, -6.0)  {Corrigir\\Submissão (RF24)};
\node[uc] (uc_dash_p)  at (7, -7.5)  {Dashboard\\Turma (RF25)};

% -- Aluno --
\node[uc] (uc_cron)    at (7, -9.0)  {Resolver c/\\Cronômetro (RF16)};
\node[uc] (uc_gatid)   at (7, -10.5) {Identificar\\Gatilhos (RF17)};
\node[uc] (uc_ferrq)   at (7, -12.0) {Ferramentas\\Qualidade (RF18--22)};
\node[uc] (uc_sub)     at (7, -13.5) {Submeter\\Resolução (RF23)};

% -- Ligações Admin --
\draw[->] (admin) -- (login);
\draw[->] (admin) -- (uc_users);
\draw[->] (admin) -- (uc_modgtt);
\draw[->] (admin) -- (uc_gat);
\draw[->] (admin) -- (uc_turmas);
\draw[->] (admin) -- (uc_dash_a);

% -- Ligações Professor --
\draw[->] (prof) -- (login);
\draw[->] (prof) -- (uc_atv);
\draw[->] (prof) -- (uc_corr);
\draw[->] (prof) -- (uc_dash_p);

% -- Ligações Aluno --
\draw[->] (aluno) -- (login);
\draw[->] (aluno) -- (uc_cron);
\draw[->] (aluno) -- (uc_gatid);
\draw[->] (aluno) -- (uc_ferrq);
\draw[->] (aluno) -- (uc_sub);

% -- Include --
\draw[->, dashed] (login) -- node[right, font=\tiny]{include} (senha);
\draw[->, dashed] (uc_cron) -- (uc_gatid);
\draw[->, dashed] (uc_gatid) -- (uc_ferrq);
\draw[->, dashed] (uc_ferrq) -- (uc_sub);

\end{tikzpicture}
\fonte{Elaborado pelo autor com TikZ (2026).}
\end{figure}

% ---------------------------------------------------------------------------
\section{Diagrama de Classes de Domínio}
\label{sec_classes_dominio}
% ---------------------------------------------------------------------------

O Diagrama de Classes de Domínio da \autoref{fig_classes} ilustra as
principais entidades do modelo de domínio (camada \textit{model/entity})
da aplicação Spring Boot, com seus atributos e as associações que orientam
o mapeamento JPA/Hibernate.

\begin{figure}[htb]
\centering
\caption{Diagrama de Classes de Domínio simplificado do SIGEA-GTT (UML 2.5).}
\label{fig_classes}

\tikzset{
  cls/.style={
    rectangle, draw=blue!70!black, fill=white, thick,
    minimum width=3.6cm, font=\scriptsize, align=left
  },
  clshead/.style={
    rectangle, draw=blue!70!black, fill=blue!15, thick,
    minimum width=3.6cm, minimum height=0.65cm,
    font=\small\bfseries, align=center
  }
}

\begin{tikzpicture}[>=Stealth, node distance=0.4cm]

% User
\node[clshead] (uh) at (0,0)   {User};
\node[cls, anchor=north] (u) at (uh.south)
  {id: UUID\\email: String\\role: UserRole\\matricula: String\\primeiroAcesso: Boolean};

% Turma
\node[clshead] (th) at (5,0)   {Turma};
\node[cls, anchor=north] (t) at (th.south)
  {id: Long\\nome: String\\status: TurmaStatus\\professor: User};

% Atividade
\node[clshead] (ah) at (10,0)  {Atividade};
\node[cls, anchor=north] (a) at (ah.south)
  {id: Long\\titulo: String\\prazo: LocalDateTime\\prontuario: JSONB};

% Submissao
\node[clshead] (sh) at (10,-4) {Submissao};
\node[cls, anchor=north] (s) at (sh.south)
  {id: Long\\aluno: User\\nota: BigDecimal\\parecer: String\\status: SubmissaoStatus};

% ModuloGTT
\node[clshead] (mh) at (0,-4)  {ModuloGTT};
\node[cls, anchor=north] (m) at (mh.south)
  {id: Long\\codigo: String\\nome: String};

% GatilhoGTT
\node[clshead] (gh) at (5,-4)  {GatilhoGTT};
\node[cls, anchor=north] (g) at (gh.south)
  {id: Long\\codigo: String\\descricao: String\\pistaClinica: String};

% GravidadeMERP
\node[clshead] (rh) at (5,-8)  {GravidadeMERP};
\node[cls, anchor=north] (r) at (rh.south)
  {id: Long\\categoria: String\\danoReal: Boolean\\descricao: String};

% Relacionamentos
\draw[->, thick] (th.west) -- node[above, font=\tiny]{1 professor} (uh.east);
\draw[->, thick] (ah.west) -- node[above, font=\tiny]{N} (th.east);
\draw[->, thick] (sh.north) -- node[right, font=\tiny]{N} (ah.south);
\draw[->, thick] (sh.west) -- node[above, font=\tiny]{N:M} (g.east);
\draw[->, thick] (gh.west) -- node[above, font=\tiny]{N} (mh.east);
\draw[->, thick] (rh.north) -- node[right, font=\tiny]{1} (gh.south);

\end{tikzpicture}
\fonte{Elaborado pelo autor com TikZ (2026).}
\end{figure}

% ---------------------------------------------------------------------------
\section{Visão Arquitetural 4+1 (Kruchten)}
\label{sec_arq_41}
% ---------------------------------------------------------------------------

O Modelo de Visão Arquitetural 4+1 \cite{kruchten1995} descreve a
arquitetura de software por meio de cinco perspectivas complementares,
cada uma endereçada a um grupo distinto de \textit{stakeholders}:
\textit{Lógica}, \textit{Processo}, \textit{Física},
\textit{Desenvolvimento} e \textit{Cenários} (a quinta vista que une
as demais). As subseções a seguir detalham cada visão no contexto do
SIGEA-GTT.

\subsection{Visão Lógica}
\label{subsec_visao_logica}

A arquitetura lógica do SIGEA-GTT segue o padrão de camadas
(\textit{Layered Architecture}) no backend Spring Boot, complementado
pela arquitetura de componentes \textit{standalone} no frontend Angular 18.
Internamente, o backend é estruturado em quatro camadas horizontais:

\begin{enumerate}
    \item \textbf{Camada de Apresentação (Controllers REST)}: Expõe os
    \textit{endpoints} JSON anotados com \texttt{@RestController} via
    SpringDoc OpenAPI 3, aplicando validação de entrada com Bean Validation
    (\texttt{@Valid}) e tratamento centralizado de exceções por
    \texttt{@ControllerAdvice};
    \item \textbf{Camada de Aplicação / Serviço}: Orquestra a lógica de
    negócio, transações, eventos de domínio e mapeamentos via MapStruct.
    Cada serviço é anotado com \texttt{@Service} e \texttt{@Transactional};
    \item \textbf{Camada de Domínio}: Entidades JPA decoradas com Lombok,
    mapeamentos de herança e relacionamentos bidirecionais; Value Objects
    para tipos complexos como \texttt{JSONB};
    \item \textbf{Camada de Infraestrutura}: Repositórios Spring Data JPA
    (\texttt{JpaRepository}), \textit{migrations} Flyway, configurações de
    segurança Spring Security 6 com JWT e provedores de e-mail.
\end{enumerate}

\subsection{Visão de Processo}
\label{subsec_visao_processo}

O fluxo de resolução de atividade pelo aluno constitui o processo mais
crítico do sistema. O \autoref{alg_fluxo_resolucao} apresenta a
sequência de estados da submissão, desde a abertura do prontuário
simulado até o parecer formativo do docente.

\begin{algoritmo}[htb]
\caption{Fluxo de Resolução de Atividade GTT (SIGEA-GTT)}
\label{alg_fluxo_resolucao}
\begin{algorithm}[H]
\KwData{Aluno autenticado; Atividade disponível com Prontuário Simulado}
\KwResult{Submissão avaliada com nota e parecer pedagógico}
Aluno acessa a atividade dentro do prazo\;
Sistema exibe prontuário fictício e inicia cronômetro de 20 min\;
\While{cronômetro > 0 e não submetido}{
  Aluno identifica gatilhos e classifica gravidade NCC MERP\;
  Aluno preenche Ishikawa, GUT, 5W2H, PDCA, SWOT e Brainstorming\;
  \If{aluno clicar em Submeter}{
    Sistema persiste submissão com timestamp e encerra cronômetro\;
    \textbf{break}\;
  }
}
\If{cronômetro = 0 e não submetido}{
  Sistema auto-submete o formulário no estado atual\;
}
Professor recebe notificação e abre tela de correção\;
Professor atribui nota (0,00--10,00) e parecer formativo\;
Sistema atualiza status para \texttt{CORRIGIDA}\;
Aluno consulta resultado e aprende com o feedback\;
\end{algorithm}
\end{algoritmo}

\subsection{Visão Física (Implantação)}
\label{subsec_visao_fisica}

A \autoref{fig_deploy} descreve a topologia de implantação do SIGEA-GTT
em ambiente universitário. O sistema pode ser executado tanto em contêineres
Podman/Docker quanto diretamente em máquinas virtuais Linux no datacenter
da UFS.

\begin{figure}[htb]
\centering
\caption{Diagrama de Implantação (Visão Física) do SIGEA-GTT.}
\label{fig_deploy}
\begin{tikzpicture}[
  >=Stealth,
  node3d/.style={
    rectangle, draw=blue!70!black, fill=blue!8, thick,
    minimum width=3.5cm, minimum height=1.0cm,
    font=\small\bfseries, rounded corners=3pt
  },
  artifact/.style={
    rectangle, draw=gray!60, fill=gray!5, thin,
    minimum width=3.0cm, minimum height=0.65cm,
    font=\scriptsize, rounded corners=1pt
  }
]
% Nós
\node[node3d] (browser) at (0, 0)   {Navegador Web};
\node[artifact] (spanode) at (0, -1.2) {Angular 18 SPA};

\node[node3d] (appsvr) at (6, 0)   {Servidor App (OpenJDK 21)};
\node[artifact] (jar)    at (6, -1.2) {SIGEA-GTT.jar};
\node[artifact] (nginx)  at (6, -2.0) {Nginx (static)};

\node[node3d] (dbsvr) at (12, 0)   {Servidor BD (PostgreSQL 16)};
\node[artifact] (pgdb)  at (12, -1.2) {sigea\_gtt (schema)};
\node[artifact] (fly)   at (12, -2.0) {Flyway V1..VN};

% Ligações
\draw[->, thick] (browser) -- node[above, font=\scriptsize]{HTTPS :443} (appsvr);
\draw[->, thick] (appsvr)  -- node[above, font=\scriptsize]{JDBC :5432} (dbsvr);

\end{tikzpicture}
\fonte{Elaborado pelo autor com TikZ (2026).}
\end{figure}

\subsection{Visão de Desenvolvimento}
\label{subsec_visao_desenv}

O repositório é organizado como um monorepo com separação clara entre
\texttt{backend/} (Spring Boot/Maven) e \texttt{frontend/} (Angular/npm).
A estrutura de pacotes Java segue a convenção
\texttt{br.ufs.sigea.<camada>.<módulo>}. As \textit{branches} de
desenvolvimento seguem o modelo Gitflow: \texttt{main} (estável),
\texttt{develop} (integração contínua) e \texttt{feature/<nome>}
(desenvolvimento isolado).

\subsection{Visão de Cenários (Casos de Uso)}
\label{subsec_visao_cenarios}

A quinta vista é o próprio Diagrama de Casos de Uso descrito na
\autoref{sec_casos_uso}, que serve de unificador narrativo das quatro
visões anteriores, confirmando que a arquitetura cobre integralmente
os requisitos RF01--RF25 e RNF01--RNF10 apresentados no
\autoref{cap_requisitos}.
